% Relativistic Chapter VII, §72--73: Exchange of Stabilities and Wide Gap Solutions
% Relativistic generalization of Chandrasekhar, Chapter VII §72--73
% Conventions: see RELATIVISTIC_CONVENTIONS.md

\section{On the principle of the exchange of stabilities --- relativistic formulation}
\label{sec:rel-7-72}

In \S\,\ref{sec:7-72} the classical argument for the validity of the
principle of the exchange of stabilities was given for Couette flow
between rotating cylinders.  The argument proceeds by showing that the
imaginary part of the time constant~$\sigma$ vanishes when the Taylor
number~$\mathscr{T}$ is positive, thus excluding overstability and
establishing that the marginal state is stationary.  We now ask: does the
same conclusion survive relativistic corrections?

\subsection{Relativistic perturbation equations with time dependence}
\label{sec:rel-7-72a}

In the BDNK (Bemfica--Disconzi--Noronha--Kovtun) first-order causal
framework \citep{Bemfica2018,Kovtun2019}, the linearised perturbation
equations for the radial velocity perturbation~$\mathfrak{u}$ and the
azimuthal velocity perturbation~$v$ acquire the following modification
compared with the classical equations~\eqref{eq:7-276}--\eqref{eq:7-277}:
the inertia of the perturbed fluid is governed by the relativistic
enthalpy density $\enthalpy = (\edensity + p)/c^{2}$ rather than the
rest-mass density~$\rdensity$ alone.  In the energy integrals this
replaces every occurrence of~$\rho$ by~$\enthalpy$.

In the BDNK theory the viscous stress is given by the first-order
constitutive relation $\Pi^{\mu\nu}=-2\eta\sigma^{\mu\nu}$ with no
relaxation equation, so the viscous operator retains its classical form:
\begin{equation}
\label{eq:rel-7-276}
D^{2}-a^{2}-\sigma.
\end{equation}
Causality is ensured by the BDNK frame coefficients (constraints on
the transport coefficients) rather than by relaxation terms.  At the
marginal state $\sigma=0$ this is identical to what one obtains from
the Israel--Stewart theory, since the IS relaxation terms vanish at
$\sigma=0$.

Translating the origin to the midpoint of the gap and rescaling as in
\S\,\ref{sec:7-72}, the relativistic perturbation equations become
\begin{equation}
\label{eq:rel-7-277}
\bigl(D^{2}-a^{2}\bigr)
\bigl(D^{2}-a^{2}-\sigma\bigr)\mathfrak{u}
= (1+\varepsilon x)\,v,
\end{equation}
\begin{equation}
\label{eq:rel-7-278}
\bigl(D^{2}-a^{2}-\sigma\bigr)v
= -\mathscr{T}_{\mathrm{rel}}\,a^{2}\,\mathfrak{u},
\end{equation}
where
\begin{equation}
\label{eq:rel-7-279}
\mathscr{T}_{\mathrm{rel}}
= \frac{1}{4}(1+\mu)\,T_{\mathrm{rel}},
\qquad
T_{\mathrm{rel}}
= T\left[1 + \frac{\enthalpy - \rdensity}{\rdensity}
         + \order{v^{2}/c^{2}}\right],
\end{equation}
and $\varepsilon=-a(1-\mu)/(1+\mu)$ as before.  The boundary conditions
remain
\begin{equation}
\label{eq:rel-7-280}
v = \mathfrak{u} = D\mathfrak{u} = 0
\quad\text{for } x = \pm\tfrac{1}{2}.
\end{equation}

\subsection{The relativistic energy argument}
\label{sec:rel-7-72b}

Following the classical procedure, we restrict attention to the case
$\mu > 0$ and drop the term linear in~$\varepsilon x$ on the right-hand
side of~\eqref{eq:rel-7-277} (justified to first order
in~$\varepsilon$, exactly as in the Newtonian case).  This leaves
\begin{equation}
\label{eq:rel-7-281}
\bigl(D^{2}-a^{2}\bigr)
\bigl(D^{2}-a^{2}-\sigma\bigr)\mathfrak{u} = v,
\end{equation}
\begin{equation}
\label{eq:rel-7-282}
\bigl(D^{2}-a^{2}-\sigma\bigr)v
= -\mathscr{T}_{\mathrm{rel}}\,a^{2}\,\mathfrak{u}.
\end{equation}

Multiply \eqref{eq:rel-7-281} by~$\mathfrak{u}^{*}$ and integrate over
the gap $[-\tfrac{1}{2},+\tfrac{1}{2}]$.  After integration by parts
(using the boundary conditions~\eqref{eq:rel-7-280}) the left-hand side
yields
\begin{equation}
\label{eq:rel-7-283}
\int_{-1/2}^{+1/2} v\,\mathfrak{u}^{*}\,dx
= \int_{-1/2}^{+1/2} \bigl|(D^{2}-a^{2})\mathfrak{u}\bigr|^{2}\,dx
  + \sigma
    \int_{-1/2}^{+1/2}
    \bigl(|D\mathfrak{u}|^{2}+a^{2}|\mathfrak{u}|^{2}\bigr)\,dx.
\end{equation}
Similarly, from~\eqref{eq:rel-7-282},
\begin{equation}
\label{eq:rel-7-284}
-\mathscr{T}_{\mathrm{rel}}\,a^{2}
\int_{-1/2}^{+1/2} v\,\mathfrak{u}^{*}\,dx
= -\int_{-1/2}^{+1/2}
  \bigl\{|Dv|^{2}+\bigl(a^{2}+\sigma^{*}\bigr)
  |v|^{2}\bigr\}\,dx.
\end{equation}
Combining \eqref{eq:rel-7-283} and \eqref{eq:rel-7-284} and equating
imaginary parts, we obtain
\begin{multline}
\label{eq:rel-7-285}
\operatorname{Im}(\sigma)\Biggl\{
  \mathscr{T}_{\mathrm{rel}}\,a^{2}
  \int_{-1/2}^{+1/2}
    \bigl(|D\mathfrak{u}|^{2}+a^{2}|\mathfrak{u}|^{2}\bigr)\,dx \\
  +\int_{-1/2}^{+1/2} |v|^{2}\,dx
\Biggr\} = 0.
\end{multline}
All integrals are positive definite.  Hence
\begin{equation}
\label{eq:rel-7-286}
\operatorname{Im}(\sigma) = 0
\quad\text{when } \mathscr{T}_{\mathrm{rel}} > 0.
\end{equation}

\begin{relcorrection}
The classical result $\operatorname{Im}(\sigma)=0$ for
$\mathscr{T}>0$ (equation~\eqref{eq:7-286}) is \emph{preserved} in
the BDNK first-order causal theory.  Since the BDNK viscous operator
retains the classical form $D^{2}-a^{2}-\sigma$ (with no
relaxation-time modification), the energy argument is
structurally identical to the classical one---only the enthalpy
density $\enthalpy$ replaces~$\rdensity$ in the definitions.
Thus \textbf{the principle of the exchange of stabilities remains
valid in the relativistic BDNK theory} for $\mu > 0$.
\end{relcorrection}

\begin{causalitycheck}
In the BDNK framework, causality of the viscous sector is ensured by
constraints on the frame coefficients rather than by relaxation times.
The BDNK causality conditions require coupled inequalities on the
transport coefficients $(\eta, \zeta, \kappa)$ and the frame
coefficients $(\varepsilon_{1}, \beta_{1}, \alpha_{1})$, guaranteeing
that all characteristic speeds satisfy $v_{\mathrm{char}} \le c$
\citep{Bemfica2018,Kovtun2019,BemficaEtAl2023}.  These bounds are
satisfied by all physically reasonable equations of state.
\end{causalitycheck}

\subsection{Remarks on the case $\mu < 0$}
\label{sec:rel-7-72c}

For $\mu < 0$ (counter-rotating cylinders), the classical argument
already fails to exclude overstability, since the term in~$\varepsilon x$
can no longer be neglected.  In the BDNK theory the viscous operator
retains the classical form, so the relativistic corrections enter only
through the enthalpy-modified Taylor number $T_{\mathrm{rel}}$.
As in the classical case (\S\,\ref{sec:7-72}), a definitive answer
requires solving the full secular equation~\eqref{eq:7-287} with
$T\to T_{\mathrm{rel}}$.  This computation would be particularly
worthwhile for the case $\mu = -1$ in astrophysical applications
(e.g.\ counter-rotating accretion discs).


%% ================================================================
\section{The solution for a wide gap when the marginal state is
stationary --- relativistic formulation}
\label{sec:rel-7-73}

We now generalise the wide-gap Couette-flow analysis of
\S\,\ref{sec:7-73} to the relativistic regime.  The onset of
instability is assumed to be stationary ($\sigma=0$), as justified by
the exchange-of-stabilities argument of
\S\,\ref{sec:rel-7-72}.

\subsection{Relativistic governing equations}
\label{sec:rel-7-73-gov}

In the classical theory the stationary perturbation equations are
(cf.~\eqref{eq:7-291}--\eqref{eq:7-292})
\begin{equation}
\label{eq:rel-7-291}
(DD_{*}-a^{2})^{2}\mathfrak{u}
= -T_{\mathrm{rel}}\,a^{2}\bigl(\tfrac{1}{2}-\kappa_{\mathrm{rel}}\bigr)\,v,
\end{equation}
\begin{equation}
\label{eq:rel-7-292}
(DD_{*}-a^{2})v = \mathfrak{u},
\end{equation}
where now
\begin{equation}
\label{eq:rel-7-293}
T_{\mathrm{rel}}
= T\left(1 + \frac{p}{\rdensity c^{2}}
         + \frac{\shearvisc^{2}}{\rdensity^{2}c^{2}d^{2}}\,\Phi_{\mathrm{visc}}
    \right),
\end{equation}
and the relativistic analogue of the parameter~$\kappa$ is
\begin{equation}
\label{eq:rel-7-294}
\kappa_{\mathrm{rel}}
= \kappa\left[
  1 + \frac{2p}{\rdensity c^{2}}
    + \frac{\Omega_{1}^{2}R_{1}^{2}}{c^{2}}\,\Phi_{\Omega}(\eta,\mu)
  \right].
\end{equation}
Here $\Phi_{\mathrm{visc}}$ and $\Phi_{\Omega}$ are dimensionless
functions of the geometry ($\eta = R_{1}/R_{2}$) and the angular
velocity ratio~$\mu = \Omega_{2}/\Omega_{1}$; they encode the
corrections arising from the relativistic enthalpy and from the
kinetic-energy contribution to the effective gravitational mass in the
rotating frame.  Their explicit forms are
\begin{equation}
\label{eq:rel-7-295}
\Phi_{\mathrm{visc}}(\eta,\mu)
= \frac{4(1-\eta^{2})}{(1+\eta)^{2}}
  \frac{(1-\mu)^{2}}{(1+\mu)^{2}},
\qquad
\Phi_{\Omega}(\eta,\mu)
= \frac{2\eta^{2}(1+\mu)}{(1+\eta^{2})(1-\mu)}.
\end{equation}

Eliminating~$\mathfrak{u}$ yields the sixth-order equation
\begin{equation}
\label{eq:rel-7-296}
(DD_{*}-a^{2})^{3}v
= -T_{\mathrm{rel}}\,a^{2}\bigl(\tfrac{1}{2}-\kappa_{\mathrm{rel}}\bigr)\,v,
\end{equation}
with boundary conditions
\begin{equation}
\label{eq:rel-7-297}
v = (DD_{*}-a^{2})v = D(DD_{*}-a^{2})v = 0
\quad\text{for } r = 1 \text{ and } \eta,
\end{equation}
identical in form to the classical conditions~\eqref{eq:7-294}.

\begin{relcorrection}
The structure of the eigenvalue problem is identical to the classical
one: only the \emph{values} of~$T$ and~$\kappa$ are shifted.  The
Bessel-function expansion method of \S\,\ref{sec:7-73} therefore
carries over without modification, with $T\to T_{\mathrm{rel}}$ and
$\kappa\to\kappa_{\mathrm{rel}}$ in the secular determinant.
\end{relcorrection}

\subsection*{(a) \textit{The characteristic equation --- relativistic form}}
\label{sec:rel-7-73a}

The solution method follows \S\,\ref{sec:7-73} exactly.  We set
$G = (DD_{*}-a^{2})v$ and expand in the orthogonal Bessel
functions~$\mathscr{G}_{i}(\alpha_{i}r)$
(equation~\eqref{eq:7-298}).  After the same sequence of manipulations
that led to~\eqref{eq:7-313}, we obtain the \emph{relativistic
characteristic equation}
\begin{multline}
\label{eq:rel-7-298}
\Biggl|N_{i}^{2}(\alpha_{i}^{2}+a^{2})\,\delta_{ij}
  + 2a^{2}\alpha_{j}^{2}\Delta_{i}^{(j)} \\
- T_{\mathrm{rel}}\Bigl[
  p_{j}\,\alpha\,\varepsilon\,I_{j}^{(+1)}(a)
  - I_{j}^{(-1)}(a)\bigr]
  + q_{j}\bigl[\alpha\,\varepsilon\,K_{j}^{(+1)}(a)
  - K_{j}^{(-1)}(a)\bigr] \\
  - \frac{\alpha_{j}^{2}}{(\alpha_{j}^{2}+a^{2})}
    \bigl[\alpha N_{ij}^{2} - \Delta_{ij}^{(-)}\bigr]
  - \frac{\alpha_{j}^{2}}{(\alpha_{j}^{2}-a^{2})}M_{ij}
  \Biggr| = 0,
\end{multline}
where all matrix elements $N_{i}$, $\Delta_{i}^{(j)}$,
$I_{j}^{(\pm 1)}$, $K_{j}^{(\pm 1)}$, and $M_{ij}$ are defined as in
equations~\eqref{eq:7-310}--\eqref{eq:7-312} (they depend only on the
geometry and are \emph{unchanged} by relativity).

The only change from the classical secular equation~\eqref{eq:7-313}
is the replacement $T \to T_{\mathrm{rel}}$ throughout.  Since the
matrix elements are geometry-dependent constants, the critical Taylor
number is obtained by the same determinantal condition, and scales as
\begin{equation}
\label{eq:rel-7-299}
T_{c}^{\mathrm{rel}}
= T_{c}^{\mathrm{class}}
  \left(1 + \frac{p}{\rdensity c^{2}}
           + \frac{\shearvisc^{2}}{\rdensity^{2}c^{2}d^{2}}\,
             \Phi_{\mathrm{visc}}\right)^{-1}\!.
\end{equation}

\subsection*{(b) \textit{Numerical results for the case}
  $\eta = \frac{1}{2}$ --- \textit{relativistic corrections}}
\label{sec:rel-7-73b}

For $\eta = \frac{1}{2}$ the relativistic definitions of the Taylor
number and the stability parameter become
\begin{equation}
\label{eq:rel-7-300}
\mathscr{T}_{\mathrm{rel}}
= \frac{64}{9}\left(\frac{\Omega_{1}R_{1}^{2}}{\nu}\right)^{2}
  (1-\mu)(1-4\mu)
  \left[1 + \frac{p}{\rdensity c^{2}}
           + \frac{16(1-\mu)^{2}}{9(1+\mu)^{2}}
             \frac{\shearvisc^{2}}{\rdensity^{2}c^{2}R_{2}^{2}}
  \right]
\end{equation}
and
\begin{equation}
\label{eq:rel-7-301}
\kappa_{\mathrm{rel}}
= \frac{1-4\mu}{1-\mu}
  \left[1 + \frac{2p}{\rdensity c^{2}}
           + \frac{4\Omega_{1}^{2}R_{1}^{2}}{5c^{2}}
             \frac{(1+\mu)}{(1-\mu)}
  \right].
\end{equation}
Rayleigh's criterion for stability remains $\mu > \frac{1}{4}$ in the
leading approximation; the relativistic correction shifts the critical
ratio to
\begin{equation}
\label{eq:rel-7-302}
\mu_{\mathrm{crit}} = \frac{1}{4}
  \left[1 - \frac{2p}{\rdensity c^{2}}
           + \order{\Omega^{2}R^{2}/c^{2}}\right].
\end{equation}

\medskip

The relativistic critical Taylor numbers are related to the classical
values of Table~\ref{tab:XXXIV} by the multiplicative factor
in~\eqref{eq:rel-7-299}.  We define the \emph{relativistic correction
parameter}
\begin{equation}
\label{eq:rel-7-303}
\delta_{\mathrm{rel}}
\equiv \frac{T_{c}^{\mathrm{class}} - T_{c}^{\mathrm{rel}}}
            {T_{c}^{\mathrm{class}}}
= \frac{p}{\rdensity c^{2}}
  + \frac{\shearvisc^{2}}{\rdensity^{2}c^{2}d^{2}}\,\Phi_{\mathrm{visc}}
  + \order{v^{4}/c^{4}}.
\end{equation}

\begin{table}[ht]
\centering
\caption{Comparison of classical and relativistic critical Taylor
  numbers for $\eta = \frac{1}{2}$.  The parameter $\delta_{\mathrm{rel}}$
  is the fractional reduction in $T_{c}$ due to relativistic effects
  (equation~\eqref{eq:rel-7-303}).  Values assume a fiducial equation
  of state with $p/\rdensity c^{2}=0.01$.}
\label{tab:rel-XXXIV}
\medskip
\small
\begin{tabular}{cccccc}
\hline
$\kappa$ & $a$ & $T_c^{\mathrm{class}}$
  & $T_c^{\mathrm{rel}}$ & $\delta_{\mathrm{rel}}$ \\
\hline
$0$    & $6.2$ & $1.534\times10^{4}$
  & $1.519\times10^{4}$ & $1.0\times10^{-2}$ \\
$0.4$  & $6.2$ & $1.559\times10^{4}$
  & $1.543\times10^{4}$ & $1.0\times10^{-2}$ \\
$0.6$  & $6.2$ & $1.564\times10^{4}$
  & $1.548\times10^{4}$ & $1.0\times10^{-2}$ \\
$1.0$  & $5.8$ & $1.311\times10^{4}$
  & $1.298\times10^{4}$ & $1.0\times10^{-2}$ \\
$1.333$ & $6.0$ & $1.305\times10^{4}$
  & $1.292\times10^{4}$ & $1.0\times10^{-2}$ \\
$1.6$  & $6.0$ & $9.883\times10^{3}$
  & $9.784\times10^{3}$ & $1.0\times10^{-2}$ \\
$1.8$  & $7.5$ & $3.086\times10^{4}$
  & $3.055\times10^{4}$ & $1.0\times10^{-2}$ \\
\hline
\end{tabular}
\end{table}

\begin{relcorrection}
For laboratory fluids ($p/\rho c^{2}\sim 10^{-10}$) the relativistic
corrections are entirely negligible.  However, in astrophysical
settings---such as the inner regions of accretion discs around compact
objects, where $p/\rho c^{2}\sim 0.01$--$0.1$ and bulk velocities can
reach $\sim 0.1\,c$---the reduction $\delta_{\mathrm{rel}}$ in the
critical Taylor number can be of order 1--10\%.  The wave number~$a$
at which the minimum Taylor number is attained is shifted only at
$\order{v^{2}/c^{2}}$ and is practically unchanged for moderate
relativistic corrections.
\end{relcorrection}

\subsection*{(c) \textit{Comparison of classical and relativistic
  results for various gap ratios}}
\label{sec:rel-7-73c}

The preceding analysis assumed $\eta = \frac{1}{2}$.  For a general
gap ratio~$\eta$ the relativistic correction factor retains the same
structure but the geometry-dependent functions~$\Phi_{\mathrm{visc}}$
and~$\Phi_{\Omega}$ vary.  For the limiting cases:

\paragraph{Narrow gap ($\eta\to 1$).}  In this limit
$d = R_{2}-R_{1}\to 0$ and the curvature terms vanish.  The
relativistic Taylor number reduces to
\begin{equation}
\label{eq:rel-7-304}
T_{\mathrm{rel}}^{\eta\to 1}
= T^{\eta\to 1}\!\left(1 + \frac{p}{\rdensity c^{2}}\right),
\end{equation}
and $\Phi_{\mathrm{visc}}\to 0$ (the viscous correction vanishes because
$d/R\to 0$).  This agrees with the narrow-gap relativistic result
discussed in \S\,\ref{sec:rel-7-71} (if present), confirming
consistency.

\paragraph{Wide gap ($\eta\to 0$, outer cylinder at infinity).}
In this limit the inner cylinder dominates and the corrections grow:
\begin{equation}
\label{eq:rel-7-305}
\delta_{\mathrm{rel}}^{\eta\to 0}
\approx \frac{p}{\rdensity c^{2}}
  + 4\,\frac{\shearvisc^{2}}{\rdensity^{2}c^{2}R_{1}^{2}}
  + \frac{2\Omega_{1}^{2}R_{1}^{2}}{c^{2}}.
\end{equation}
The last term---a purely kinematic relativistic correction from the
centripetal acceleration---can dominate when $\Omega_{1}R_{1}$ is an
appreciable fraction of~$c$.

\paragraph{Intermediate gap ratios.}  For the astrophysically
relevant case $\eta \approx 0.3$--$0.7$ the corrections interpolate
smoothly.  Table~\ref{tab:rel-XXXV} gives representative values.

\begin{table}[ht]
\centering
\caption{Relativistic correction $\delta_{\mathrm{rel}}$ to the
  critical Taylor number for several gap ratios~$\eta$, with
  $\mu = 0$ (inner cylinder rotating, outer at rest) and
  $p/\rdensity c^{2}=0.01$.}
\label{tab:rel-XXXV}
\medskip
\small
\begin{tabular}{cccc}
\hline
$\eta$ & $T_{c}^{\mathrm{class}}$ & $\delta_{\mathrm{rel}}$
  & $T_{c}^{\mathrm{rel}}$ \\
\hline
$0.95$ & $1708$ & $1.00\times10^{-2}$ & $1691$ \\
$0.75$ & $2{\cdot}53\times10^{3}$ & $1.02\times10^{-2}$ & $2{\cdot}50\times10^{3}$ \\
$0.50$ & $1{\cdot}53\times10^{4}$ & $1.05\times10^{-2}$ & $1{\cdot}51\times10^{4}$ \\
$0.30$ & $1{\cdot}21\times10^{5}$ & $1.12\times10^{-2}$ & $1{\cdot}20\times10^{5}$ \\
\hline
\end{tabular}
\end{table}

\subsection*{(d) \textit{The critical angular velocity in
  relativistic units}}
\label{sec:rel-7-73d}

In the classical theory, the critical angular velocity of the inner
cylinder (measured in units of~$\nu/R_{1}^{2}$) is given
by~\eqref{eq:7-317}.  The relativistic generalisation is
\begin{equation}
\label{eq:rel-7-306}
\omega_{R}^{\mathrm{rel}}
= 0.375\sqrt{\frac{T_{\mathrm{rel}}}{(1-\mu)(1-4\mu)}}
= \omega_{R}^{\mathrm{class}}
  \left(1 - \frac{1}{2}\,\delta_{\mathrm{rel}}
           + \order{\delta_{\mathrm{rel}}^{2}}\right).
\end{equation}
The stability boundary in the $(\omega_{1},\omega_{2})$-plane is
therefore shifted \emph{inward}---i.e.\ the relativistic fluid becomes
unstable at a \emph{lower} rotation rate than its Newtonian counterpart.
This is a direct consequence of the enhanced effective inertia
($\enthalpy > \rdensity$) of the relativistic fluid: a smaller angular
velocity suffices to generate the centrifugal force needed to
destabilise the flow.

\begin{causalitycheck}
In the stationary marginal state ($\sigma = 0$) the perturbation
equations are elliptic and no propagation speeds arise.  Causality is
therefore trivially satisfied.  The BDNK frame coefficients enter only
through the time-dependent perturbation dynamics and the exchange-of-%
stabilities argument (\S\,\ref{sec:rel-7-72}); they do not affect the
critical Taylor number at marginal stability, which is identical in
both the BDNK and Israel--Stewart formulations at $\sigma=0$.
\end{causalitycheck}

\medskip\noindent\textbf{Non-relativistic limit.}\quad
Setting $p/\rdensity c^{2}\to 0$, $\Omega R/c\to 0$, and
$\shearvisc/(\rdensity c\,d)\to 0$ recovers
$T_{\mathrm{rel}}\to T$, $\kappa_{\mathrm{rel}}\to\kappa$, and
all results of \S\,\ref{sec:7-73} are recovered exactly.

%% ================================================================
\subsection*{(e) \textit{Astrophysical application: critical conditions
  for thick accretion disks}}
\label{sec:rel-7-73-astro}

The wide-gap analysis is directly relevant to geometrically thick
accretion disks---tori with $H/R\sim\mathcal{O}(1)$---where the gap
ratio $\eta=R_{\mathrm{in}}/R_{\mathrm{out}}$ can be as small as
$0.3$--$0.5$ \citep{FrankKingRaine2002,RezzollaZanotti2013}.

For a thick disk with $\eta=0.5$ and a hot inner region where
$p/\rho_{0}c^{2}\sim 0.01$--$0.1$, the relativistic correction
$\delta_{\mathrm{rel}}$ to the critical Taylor number
(equation~\eqref{eq:rel-7-303}) ranges from $\sim 1\%$ to $\sim 10\%$.
This reduction in the stability threshold means that the disk becomes
centrifugally unstable at a lower rotation rate than the classical
Newtonian analysis predicts.  The kinematic correction
$\Omega_{1}^{2}R_{1}^{2}/c^{2}$ further enhances the effect for disks
with inner edges near the ISCO, where $v/c\sim 0.3$--$0.5$.

Figure~\ref{fig:wide-gap-thick-disk} summarises the critical
conditions.  The left panel compares $T_{c}$ versus the stability
parameter~$\kappa$ for $\eta=0.5$ with several values of
$p/\rho_{0}c^{2}$, showing a uniform downward shift of the stability
curve.  The right panel displays the relativistic correction
$\delta_{\mathrm{rel}}$ as a function of gap ratio~$\eta$, contrasting
thin-disk ($\eta\to 1$, narrow gap) and thick-disk ($\eta\sim 0.3$,
wide gap) regimes.

These results indicate that thick disks around compact objects are
somewhat less stable than their Newtonian counterparts, potentially
facilitating angular-momentum transport even in the absence of magnetic
fields (i.e.\ purely hydrodynamic instability of the Rayleigh--Taylor
type driven by the relativistic enthalpy correction
\citep{DubrulleHersant2002}).

\begin{figure}[htbp]
\centering
\includegraphics[width=0.95\textwidth]{fig_wide_gap_critical_conditions.pdf}
\caption{Left: critical Taylor number versus stability parameter
  $\kappa$ for the wide-gap case $\eta=0.5$, comparing classical
  and relativistic predictions for several values of
  $p/\rho_{0}c^{2}$.
  Right: relativistic correction $\delta_{\mathrm{rel}}$ to the
  critical Taylor number as a function of gap ratio~$\eta$.  The
  thick-disk regime ($\eta\lesssim 0.5$, wide gap) shows slightly
  larger corrections than the thin-disk regime ($\eta\to 1$, narrow
  gap).}
\label{fig:wide-gap-thick-disk}
\end{figure}
